\documentclass{article}

\usepackage{packages}
\usepackage{environments}
\usepackage{commands}
\usepackage{titlepage}
\usepackage{svg}

\svgpath{{./}}

\setUDK{004.032.26}
\setToProgram
\setTitle{Нейронные сети с нуля}

%\setStageOne
%\setStageTwo
\setStageFinal

\setGroup{245}
\setStudentSgn{подпись студента}
\setStudent{Аветисов Тимофей Михайлович}
\setStudentDate{24.05.2026}
\setAdvisor{Трушин Дмитрий Витальевич, к.ф.-м.н.}
\setAdvisorTitle{Доцент}
\setAdvisorAffiliation{ФКН НИУ ВШЭ}
\setAdvisorDate{}
\setGrade{}
\setAdvisorSgn{подпись руководителя}
\setYear{2026}

\setmainfont{Times New Roman}
\begin{document}

\makeTitlePage
\tableofcontents

\begin{abstract}
Проект посвящен разработке библиотеки \texttt{NNS} для построения и
обучения полносвязных нейронных сетей на языке C++23 без использования
готовых фреймворков машинного обучения. В отчете рассматриваются
математическая модель прямого и обратного прохода, архитектура
программной реализации, организация обобщенных интерфейсов через
стирание типов, реализация функций потерь, оптимизаторов, загрузки
данных и обучающего цикла. Приводится пример обучения на датасете
\texttt{MNIST}.
\end{abstract}

\section{Введение}

Нейронные сети являются одним из основных инструментов современного
машинного обучения. Их применяют в задачах классификации, регрессии,
распознавания изображений, обработки сигналов, анализа естественного
языка и в других областях, где зависимость между входными признаками и
целевой переменной имеет сложный нелинейный характер. Практическое
распространение нейронных сетей связано не только с развитием
вычислительных ресурсов, но и с появлением удобных библиотек,
скрывающих детали линейной алгебры, автоматического дифференцирования,
оптимизации и обработки данных.

Высокоуровневые фреймворки позволяют быстро строить модели, однако при
их использовании часть существенных технических деталей остается
скрытой за готовыми абстракциями. В частности, пользователь часто не
работает напрямую с формами матриц, порядком вычисления градиентов,
устройством кэшей прямого прохода, состоянием оптимизатора и численной
устойчивостью функций потерь. В данном проекте эти механизмы
реализованы явно: зафиксированы соглашения о представлении данных,
выведены формулы обратного распространения ошибки и спроектированы
интерфейсы, которые допускают расширение системы новыми компонентами.

Цель проекта состоит в разработке header-only библиотеки \texttt{NNS},
предназначенной для построения и обучения полносвязных нейронных сетей
на C++23. Под выражением <<с нуля>> в рамках работы понимается отказ
от специализированных библиотек машинного обучения: проект не
использует PyTorch, TensorFlow, Keras и аналогичные средства. При этом
используются библиотеки общего назначения: Eigen для матричных
операций, \texttt{proxy} для стирания типов и Catch2 для
автоматического тестирования.

В результате работы реализована header-only библиотека, в которой
пользователь может задать последовательность линейных слоев и
скалярных функций активации, выбрать функцию потерь, оптимизатор и
стратегию изменения шага обучения, загрузить данные из матриц или
CSV-файла, а затем запустить обучение.
Архитектура построена вокруг строгого разделения ролей: слой вычисляет
локальное преобразование и локальный градиент, нейронная сеть
организует прямой и обратный проход по слоям, оптимизатор обновляет
параметры, а тренер связывает модель, данные, функцию потерь и
оптимизатор в единый цикл обучения.

В отчете сначала формулируются требования к программному проекту,
затем вводится математическая модель полносвязной сети и ее обучения.
Далее описывается программная архитектура библиотеки, внутренние
соглашения о представлении данных, устройство слоев, функций потерь,
оптимизаторов, загрузчика данных и тренера. После этого
рассматриваются пример использования на задаче классификации
изображений MNIST, организация сборки и тестирования.

\section{Функциональные и нефункциональные требования}

\subsection{Назначение библиотеки}

Библиотека \texttt{NNS} предназначена для создания, обучения и
тестового применения полносвязных нейронных сетей. Основным
пользовательским сценарием является следующий процесс: пользователь
формирует обучающую выборку в виде матриц признаков и целевых
переменных, создает объект нейронной сети из набора слоев, выбирает
функцию потерь и оптимизатор, передает эти объекты в тренер и получает
историю значения функции потерь по эпохам.

Система проектируется как библиотека, а не как отдельное приложение.
Это означает, что основная ценность проекта заключается в
предоставляемом API: классы и функции должны быть пригодны для
подключения к стороннему CMake-проекту, для написания собственных
экспериментов и для расширения пользовательскими компонентами. Файл
\texttt{main.cpp} в репозитории является примером применения
библиотеки, а не обязательной частью ее публичного интерфейса.

\subsection{Функциональные требования}

К функциональным требованиям проекта относятся:
\begin{enumerate}
    \item поддержка базовых числовых типов и матричного представления
        данных; 
    \item создание линейных слоев с настраиваемыми
        размерами входа и выхода;
    \item инициализация весов различными распределениями;
    \item поддержка встроенных функций активации \texttt{ReLU},
        \texttt{Sigmoid} и \texttt{Tanh};
    \item возможность добавления пользовательских скалярных функций
        активации;
    \item возможность выполнять предсказание, прямой проход и
        обратный проход в процессе обучения;
    \item поддержка нескольких функций потерь для задач регрессии и
        классификации;
    \item поддержка оптимизаторов \texttt{SGD} и \texttt{Adam};
    \item поддержка постоянного и убывающего шага обучения;
    \item загрузка данных из CSV-файла с выделением целевых колонок;
    \item разбиение данных на батчи с опциональным перемешиванием;
    \item запуск обучения на заданное число эпох;
    \item возврат истории среднего значения функции потерь по эпохам;
    \item наличие автоматических тестов для основных компонентов.
\end{enumerate}

Отдельным требованием является расширяемость. Пользователь должен
иметь возможность определить собственную функцию потерь, оптимизатор,
scheduler шага обучения или функцию активации при условии, что
пользовательский тип реализует ожидаемый набор методов. Для этого в
проекте используются type-erased обертки на основе библиотеки
\texttt{proxy}.

\subsection{Нефункциональные требования}

К нефункциональным требованиям относятся переносимость,
воспроизводимость, простота сборки и явная диагностика ошибок. Проект
написан на C++23 и собирается через CMake. Библиотека реализована как
\texttt{INTERFACE}-цель, поэтому ее заголовки могут подключаться без
отдельного этапа компиляции статической или динамической библиотеки.

Воспроизводимость обеспечивается классом \texttt{RandomGenerator},
который хранит генератор псевдослучайных чисел и позволяет явно
задавать seed. Это важно как для инициализации весов, так и для
перемешивания данных в \texttt{DataLoader}. При одинаковом seed и
одинаковом порядке вызовов пользователь должен получать одинаковую
последовательность случайных величин.

Диагностика ошибок реализуется через исключения стандартной
библиотеки. Компоненты проверяют размеры матриц, непустоту входов,
корректность гиперпараметров и согласованность кэшей. Такой подход
повышает надежность библиотеки: ошибка формы данных обнаруживается в
точке нарушения контракта, а не приводит к неявно неверному результату
в конце обучения.

\section{Математическая постановка}

\subsection{Общий принцип обучения и устройства нейронной сети}

Рассматривается задача обучения с учителем. Дана выборка
\[
    \mathcal{D} =
    \{(x^{(j)}, y^{(j)})\}_{j=1}^{N},
\]
где $x^{(j)} \in \mathbb{R}^{d}$ --- вектор признаков, а
$y^{(j)}$ --- целевая переменная. Требуется построить параметрическое
отображение
\[
    f_{\theta}: \mathbb{R}^{d} \to \mathbb{R}^{k},
\]
которое минимизирует среднюю ошибку на обучающей выборке.

Полносвязная нейронная сеть задается как композиция преобразований
\[
    f_{\theta}
    =
    f_m \circ f_{m-1} \circ \ldots \circ f_1.
\]
Каждое преобразование называется слоем. В полносвязной сети
параметризованные слои являются аффинными преобразованиями, а
нелинейность вносится функциями активации. Если обозначить
$h_0 = x$, то прямой проход имеет вид
\[
    h_i = f_i(h_{i-1}), \qquad i = 1,\ldots,m.
\]
Выход $h_m$ является предсказанием модели.

Обучение состоит в минимизации эмпирического риска
\[
    J(\theta)
    =
    \frac{1}{N}
    \sum_{j=1}^{N}
    \ell(f_{\theta}(x^{(j)}), y^{(j)}),
\]
где $\ell$ --- функция потерь. На каждой итерации сначала
вычисляется предсказание, затем значение функции потерь, после чего
по правилу дифференцирования композиции находятся градиенты по
параметрам. Далее параметры изменяются в направлении уменьшения
эмпирического риска.

\subsection{Представление данных}

Батч --- это поднабор объектов
$x^{(i_1)}, x^{(i_2)}, \dots, x^{(i_B)}$ и соответствующих им
целевых переменных
$y^{(i_1)}, y^{(i_2)}, \dots, y^{(i_B)}$. Для батча из $B$
объектов удобно объединить входы в матрицу
\[
    X =
    \begin{bmatrix}
        x^{(1)} & x^{(2)} & \ldots & x^{(B)}
    \end{bmatrix}
    \in \mathbb{R}^{d \times B}.
\]
Каждый столбец соответствует одному объекту. Если выходная размерность
равна $k$, то целевые значения батча можно записать как
\[
    Y =
    \begin{bmatrix}
        y^{(1)} & y^{(2)} & \ldots & y^{(B)}
    \end{bmatrix}
    \in \mathbb{R}^{k \times B}.
\]
В задачах многоклассовой классификации $y^{(j)}$ часто является
one-hot вектором класса: вектором длины, равной числу классов, с
нулями во всех позициях, кроме позиции истинного класса. Также целевое
значение может задаваться меткой класса. В задачах регрессии
$y^{(j)}$ содержит одно или несколько вещественных целевых значений.

\subsection{Линейный слой}

Пусть вход батча имеет размерность $X \in
\mathbb{R}^{d_{in} \times B}$. Линейный слой с параметрами
\[
    A \in \mathbb{R}^{d_{out} \times d_{in}},
    \qquad
    b \in \mathbb{R}^{d_{out}}
\]
задает преобразование
\[
    Z = AX + b\mathbf{1}^{T},
\]
где $\mathbf{1} \in \mathbb{R}^{B}$ --- вектор из единиц. Слагаемое
$b\mathbf{1}^{T}$ означает, что один и тот же вектор смещения
прибавляется ко всем столбцам батча.

Пусть известен градиент функции потерь по выходу слоя:
\[
    G_Z =
    \frac{\partial L}{\partial Z}
    \in \mathbb{R}^{d_{out} \times B}.
\]
Тогда градиенты по параметрам и входу выражаются формулами
\[
    \frac{\partial L}{\partial A} = G_Z X^T,
    \qquad
    \frac{\partial L}{\partial b} = G_Z \mathbf{1},
    \qquad
    \frac{\partial L}{\partial X} = A^T G_Z.
\]
Эти формулы следуют из дифференциала
$dZ = dA \cdot X + A \cdot dX + db \cdot \mathbf{1}^{T}$.

\subsection{Функции активации}

Если сеть состоит только из аффинных преобразований, то вся композиция
также является аффинным преобразованием. Поэтому между
параметризованными слоями вводятся нелинейные функции активации.
Обычно они применяются поэлементно:
\[
    H_{ij} = \varphi(Z_{ij}).
\]

В качестве примеров используются следующие функции. ReLU задается как
\[
    \operatorname{ReLU}(x) = \max(0, x),
\]
а ее производная почти всюду равна
\[
    \operatorname{ReLU}'(x) =
    \begin{cases}
        1, & x > 0,\\
        0, & x \leq 0.
    \end{cases}
\]
Сигмоида и ее производная имеют вид
\[
    \sigma(x) =
    \frac{1}{1 + e^{-x}},
    \qquad
    \sigma'(x) =
    \sigma(x)(1-\sigma(x)).
\]
Гиперболический тангенс задается формулой
\[
    \tanh(x) =
    \frac{e^x - e^{-x}}{e^x + e^{-x}},
\]
а его производная равна
\[
    \tanh'(x) = 1 - \tanh^2(x).
\]
Эти функции подробно рассматриваются в литературе по
глубокому обучению~\cite{goodfellow2016deep}.

\subsection{Функции потерь}

Функция потерь $\ell(\hat{y}, y)$ измеряет расхождение между
предсказанием $\hat{y}$ и целевым значением $y$.

Для регрессии часто используется среднеквадратичная ошибка:
\[
    L_{MSE}(\hat{Y}, Y) =
    \frac{1}{n}
    \sum_{i,j}
    (\hat{Y}_{ij} - Y_{ij})^2,
\]
где $n$ --- число элементов матрицы. Ее градиент по предсказанию:
\[
    \frac{\partial L_{MSE}}{\partial \hat{Y}}
    =
    \frac{2}{n}(\hat{Y} - Y).
\]
Средняя абсолютная ошибка задается формулой
\[
    L_{MAE}(\hat{Y}, Y) =
    \frac{1}{n}
    \sum_{i,j}
    |\hat{Y}_{ij} - Y_{ij}|.
\]
Функция Хьюбера для ошибки $a = \hat{y} - y$ определяется как
\[
    L_{\delta}(a) =
    \begin{cases}
        \frac{1}{2}a^2, & |a| \leq \delta,\\
        \delta(|a| - \frac{1}{2}\delta), & |a| > \delta.
    \end{cases}
\]
Она ведет себя как квадратичная ошибка при малых отклонениях и как
абсолютная ошибка при больших отклонениях.

Для бинарной классификации используется бинарная кросс-энтропия
\[
    L_{BCE}(p, y)
    =
    -y\log p
    -(1-y)\log(1-p),
\]
где $p \in (0,1)$ --- предсказанная вероятность класса $1$.
Если модель возвращает logit $z$, то $p=\sigma(z)$. В этом случае
кросс-энтропию можно записать напрямую через $z$:
\[
    L(z, y)
    =
    \max(z, 0) - zy + \log(1 + e^{-|z|}).
\]
Такая запись численно устойчивее прямого вычисления
$-\log \sigma(z)$ или $-\log(1-\sigma(z))$, поскольку не требует
вычислять логарифм вероятности, которая из-за округления может стать
слишком близкой к нулю.

\subsection{Softmax и кросс-энтропия}

В многоклассовой классификации модель возвращает вектор logits
$z \in \mathbb{R}^{k}$. Чтобы получить распределение вероятностей по
классам, применяется softmax:
\[
    p_i =
    \frac{e^{z_i}}{\sum_{r=1}^{k} e^{z_r}},
    \qquad i=1,\ldots,k.
\]
Для one-hot вектора $y$ кросс-энтропия равна
\[
    L(z, y)
    =
    -\sum_{i=1}^{k}
    y_i \log p_i.
\]
Если истинный класс имеет номер $c$, то это выражение сводится к
\[
    L(z, y) = -\log p_c.
\]

На практике softmax и кросс-энтропию целесообразно рассматривать как
единую функцию потерь, принимающую logits. Причина состоит в численной
устойчивости. Экспоненты $e^{z_i}$ могут переполниться при больших
положительных $z_i$. Поэтому используется тождественное преобразование
\[
    \operatorname{softmax}(z)
    =
    \operatorname{softmax}(z - \alpha\mathbf{1})
\]
для любой константы $\alpha$. Обычно выбирают
\[
    \alpha = \max_i z_i.
\]
Тогда все аргументы экспоненты неположительны, а максимальный равен
нулю:
\[
    p_i =
    \frac{e^{z_i-\alpha}}
         {\sum_{r=1}^{k} e^{z_r-\alpha}}.
\]

Кроме того, производная объединенной функции имеет простой вид. Для
одного объекта
\[
    \frac{\partial L}{\partial z_i}
    =
    p_i - y_i.
\]
Для батча размера $B$ при усреднении по объектам градиент делится на
$B$. Поэтому хранить softmax как отдельный слой перед
кросс-энтропией не обязательно: математически корректнее и численно
устойчивее объединить эти две операции в одну функцию потерь,
работающую с logits.

\subsection{Обратное распространение ошибки}

Обратное распространение ошибки является применением правила
дифференцирования композиции функций к последовательности слоев.
Подробный вывод алгоритма backpropagation для полносвязных сетей
приведен, например, в книге M.
Nielsen~\cite{nielsen2015neural}.

Пусть
\[
    h_i = f_i(h_{i-1}; \theta_i),
    \qquad i=1,\ldots,m,
\]
а итоговая функция потерь равна
\[
    L = \ell(h_m, y).
\]
Обозначим
\[
    G_i =
    \frac{\partial L}{\partial h_i}.
\]
Тогда обратный проход вычисляет градиенты в порядке
$i=m,m-1,\ldots,1$:
\[
    G_{i-1}
    =
    \left(
        \frac{\partial h_i}{\partial h_{i-1}}
    \right)^T
    G_i.
\]
Если слой имеет параметры $\theta_i$, то одновременно вычисляется
\[
    \frac{\partial L}{\partial \theta_i}
    =
    \left(
        \frac{\partial h_i}{\partial \theta_i}
    \right)^T
    G_i.
\]

После вычисления градиентов параметры обновляются оптимизационным
методом. В простейшем случае градиентного спуска
\[
    \theta_{t+1}
    =
    \theta_t
    -
    \eta_t
    \nabla_{\theta} J(\theta_t),
\]
где $\eta_t > 0$ --- шаг обучения. Более сложные методы, например
Adam, используют дополнительные оценки моментов градиента, но также
основаны на градиенте эмпирического риска.

\section{Общая архитектура проекта}

\subsection{Структура директорий}

Публичные заголовки библиотеки расположены в директории
\texttt{include/nns}. Тесты находятся в папке \texttt{tests}.
Документация расположена в \texttt{docs}, а отчет --- в
\texttt{docs/formal}. Главный
публичный заголовок \texttt{include/nns/NNS.hpp} подключает все
основные модули библиотеки и предназначен для пользовательского кода.

\newpage
Основные группы заголовков:
\begin{itemize}
    \item \texttt{core} --- базовые типы и strong type-обертки;
    \item \texttt{activation} --- функции активации и type-erased
        интерфейс активации;
    \item \texttt{layers} --- линейный слой,
        слой активации и type-erased интерфейс слоя;
    \item \texttt{network} --- класс нейронной сети;
    \item \texttt{loss} --- функции потерь и их обобщенный интерфейс; 
    \item \texttt{learningrates} --- стратегии изменения шага
        обучения;
    \item \texttt{optimizer} --- оптимизаторы и их обобщенный
        интерфейс; 
    \item \texttt{trainer} --- цикл обучения; 
    \item \texttt{utils} --- генератор случайных чисел,
        загрузчик данных и CSV-парсер.
\end{itemize}

На рисунке~\ref{fig:components} приведена компонентная схема
библиотеки. Центральным объектом в процессе обучения является
\texttt{Trainer}: он получает батчи от \texttt{DataLoader}, вызывает
прямой и обратный проход \texttt{NeuralNetwork}, использует
\texttt{Loss} для вычисления значения ошибки и начального градиента, а
затем передает градиенты в \texttt{Optimizer}.

\begin{figure}[!t]
    \centering
    \includesvg[
        inkscapelatex=false,
        width=1\linewidth
    ]{component_architecture}
    \caption{Компонентная архитектура библиотеки \texttt{NNS}}
    \label{fig:components}
\end{figure}

\subsection{Типы данных}
В коде используются базовые типы, заданные в
\texttt{include/nns/core/Types.hpp}:
\begin{cppcode}
using Scalar = double;
using Matrix = Eigen::Matrix<Scalar, Eigen::Dynamic, Eigen::Dynamic>;
using Vector = Eigen::Matrix<Scalar, Eigen::Dynamic, 1>;
using Index = Eigen::Index;
\end{cppcode}
Тип \texttt{Scalar} задает числовой тип библиотеки.
\texttt{Matrix} используется для батчей, весов и градиентов весов.
\texttt{Vector} используется для смещений линейных слоев и их
градиентов.
\subsection{Внешние зависимости}

Для матричных операций используется Eigen версии
3.4.0~\cite{eigenweb}. Eigen предоставляет динамические матрицы,
векторизованные операции, поэлементные выражения и удобные средства
для работы с блоками матриц. Это позволяет реализовать вычисления
нейронной сети без ручного управления памятью для двумерных массивов.

Для стирания типов используется библиотека \texttt{proxy} версии
4.0.0~\cite{proxygithub}. Она позволяет описывать фасады с требуемыми
методами и хранить объекты разных типов за единым
интерфейсом без классического наследования с виртуальными методами.
В проекте этот механизм применяется для активаций, слоев, функций
потерь, оптимизаторов и scheduler-ов.

Для тестирования используется Catch2 версии 3.7.1~\cite{catch2github}.
Тесты собираются как отдельный исполняемый файл \texttt{nns\_tests} и
регистрируются в CTest через \texttt{catch\_discover\_tests}.
Зависимости подтягиваются через CMake
\texttt{FetchContent}~\cite{cmakefetchcontent}.


\section{Классы модели и обобщенные интерфейсы}

\subsection{Основные классы библиотеки}

Публичный API библиотеки состоит из нескольких групп классов. Их роли
разделены так, чтобы модель, данные, функция потерь и процесс обучения
могли тестироваться и заменяться независимо.

\begin{itemize}
    \item \texttt{RandomGenerator} --- генератор случайных чисел,
        перемешивание индексов и инициализация матриц;
    \item \texttt{LinearLayer} --- Линейный слой;
    \item \texttt{ActivationLayer} --- Слой применяющий функцию активации;
    \item \texttt{NeuralNetwork} --- Нейронная сеть;
    \item \texttt{MSELoss}, \texttt{MAELoss}, \texttt{HuberLoss},
        \texttt{BCELoss}, \texttt{CrossEntropyLoss} ---
        функции потерь;
    \item \texttt{ConstantLearningRate} и
        \texttt{ExponentialLearningRate} --- scheduler-ы шага
        обучения;
    \item \texttt{SGDOptimizer} и \texttt{AdamOptimizer} ---
        оптимизаторы параметров;
    \item \texttt{DataLoader} --- разбиение матриц данных на батчи;
    \item \texttt{Trainer} --- координатор цикла обучения.
\end{itemize}

Для расширяемых точек API используются обобщенные оболочки:
\texttt{AnyScalarActivation}, \texttt{AnyLayer},

\texttt{AnyLossFunction}, \texttt{AnyOptimizer} и
\texttt{AnyLearningRateScheduler}. Они позволяют хранить
пользовательские типы за единым интерфейсом без общего
базового класса.

\subsection{Классы слоев и сети}

Класс \texttt{LinearLayer} хранит матрицу весов \texttt{A\_} и вектор смещения \texttt{b\_}.
Конструктор принимает входную и выходную размерность, так же необязательные параметры
генератор случайных чисел, схему распределения(по умлочанию нормальное распределение) и коэффициент масштаба(по умлочанию 1):
\begin{cppcode}
nns::LinearLayer layer(
    nns::In{784},
    nns::Out{128},
    rng,
    nns::Distribution::HeNormal,
    nns::Gain{1.0});
\end{cppcode}
При создании слоя проверяется, что обе размерности положительны. Затем
создается матрица весов размера \texttt{out x in}, а вектор смещения
инициализируется нулями. Нулевая инициализация смещения не создает
симметрию между нейронами, поскольку симметрия уже нарушается
случайной инициализацией весов.

Класс \texttt{ActivationLayer} хранит \texttt{AnyScalarActivation}.
Его задача --- применить скалярную функцию активации к каждому
элементу матрицы. Такой слой не имеет обучаемых параметров, но
участвует в обратном проходе: входящий градиент умножается на
производную функции активации в точке, сохраненной при прямом проходе.

Класс \texttt{NeuralNetwork} хранит последовательность
\texttt{AnyLayer}. Конструктор является variadic template и
принимает произвольную последовательность объектов, которые можно
преобразовать в слой:
\begin{cppcode}
nns::NeuralNetwork net(
    nns::LinearLayer(
        nns::In{784},
        nns::Out{128},
        rng,
        nns::Distribution::HeNormal),
    nns::ReLU(),
    nns::LinearLayer(
        nns::In{128},
        nns::Out{10},
        rng,
        nns::Distribution::XavierNormal));
\end{cppcode}
Линейные слои передаются как готовые объекты, а скалярные функции
активации автоматически оборачиваются в \texttt{ActivationLayer}. В
результате пользовательский код остается близким к математическому
описанию сети: линейное преобразование, нелинейность, следующее
линейное преобразование.

\subsection{Стирание типов: контракты компонентов}

Нейронная сеть должна хранить слои разных типов в одном контейнере.
Такая же задача возникает для функций потерь, оптимизаторов и
scheduler-ов. В проекте для этого используется стирание типов через
\texttt{proxy}. Конкретный тип объекта скрывается, но требуемый набор
методов фиксируется фасадом(интерфейсом).

Для скалярной функции активации требуются два метода:
\begin{cppcode}
Scalar forward(Scalar x) const;
Scalar derivative(Scalar x) const;
\end{cppcode}
Фасад \texttt{ScalarActivation} описывает эти операции, а тип
\texttt{AnyScalarActivation} хранит объект любой функции активации,
удовлетворяющей этому интерфейсу. Например, пользователь может
определить собственную функцию \texttt{LeakyReLU}:
\begin{cppcode}
struct LeakyReLU {
    nns::Scalar forward(nns::Scalar x) const {
        return x > 0.0 ? x : 0.01 * x;
    }

    nns::Scalar derivative(nns::Scalar x) const {
        return x > 0.0 ? 1.0 : 0.01;
    }
};
\end{cppcode}

Слой должен поддерживать четыре операции:
\begin{cppcode}
std::pair<Matrix, std::any> forward(Matrix&& X);
Matrix predict(const Matrix& X) const;
std::pair<Matrix, std::any> backward(
    Matrix&& dY,
    const std::any& cache);
std::any update(
    std::any&& grads,
    AnyOptimizer& opt,
    std::any&& opt_cache);
\end{cppcode}
Метод \texttt{predict} используется для предсказания и не создает и не сохраняет кэш.
Метод \texttt{forward} используется при обучении и возвращает кэш для
последующего обратного прохода. Метод \texttt{backward} принимает
градиент по выходу слоя и кэш прямого прохода, а возвращает градиент
по входу и градиенты параметров. Метод \texttt{update} применяет
оптимизатор к параметрам слоя и возвращает новое состояние
оптимизатора для этого слоя.

Функция потерь предоставляет два метода:
\begin{cppcode}
Scalar loss(const Matrix& y_hat, const Matrix& y);
Matrix gradient(const Matrix& y_hat, const Matrix& y);
\end{cppcode}
\newpage
Оптимизатор обновляет матричные и векторные параметры:
\begin{cppcode}
std::any update_weights(
    Matrix& param,
    const Matrix& grad,
    std::any&& cache);
std::any update_weights(
    Vector& param,
    const Vector& grad,
    std::any&& cache);
void iter_step();
\end{cppcode}
Scheduler шага обучения предоставляет текущий шаг, увеличивает номер
итерации и возвращает номер текущей итерации:
\begin{cppcode}
Scalar get_lr();
void iter_step();
size_t get_iter() const;
\end{cppcode}
Такое разделение не смешивает расписание шага обучения с правилом
обновления параметров. Один и тот же scheduler может использоваться в
разных оптимизаторах.

\subsection{Поток данных при обучении}

Обучение одного батча строится как последовательность трех проходов по
сети. Сначала \texttt{NeuralNetwork::forward} передает входную матрицу
по слоям слева направо. Каждый слой возвращает выходную матрицу и свой
локальный кэш:
\begin{cppcode}
auto [Y, cache] = layer->forward(std::move(X));
X = std::move(Y);
layers_cache.push_back(std::move(cache));
\end{cppcode}
Сеть не интерпретирует содержимое локального кэша. Она только
сохраняет кэши в \texttt{std::vector<std::any>} в том же порядке, в
котором расположены слои.

\begin{figure}[!t]
    \centering
    \includesvg[inkscapelatex=false,width=1\linewidth]{forward_pass}
    \caption{Прямой проход по последовательности слоев}
    \label{fig:forward}
\end{figure}

После вычисления значения функции потерь вызывается
\texttt{Loss::gradient}. Полученный градиент является начальным
градиентом для \texttt{NeuralNetwork::backward}. Обратный проход идет
по слоям справа налево. Для слоя с индексом \texttt{i} сеть берет
кэш с тем же индексом, который был создан этим слоем в прямом проходе.
Таким образом, слой получает ровно тот локальный объект состояния,
который он сам вернул из \texttt{forward}.

Линейный слой сохраняет входную матрицу батча, потому что она нужна
для вычисления градиента по весам. Слой активации также сохраняет
вход, так как производная активации вычисляется в точке прямого
прохода. При этом типы кэшей у разных слоев не обязаны совпадать:
конкретный слой сам знает, какой тип он положил в \texttt{std::any},
и сам извлекает его в \texttt{backward}.

\begin{figure}[!t]
    \centering
    \includesvg[
        inkscapelatex=false,
        width=1\linewidth
    ]{backward_update}
    \caption{Обратный проход и обновление параметров}
    \label{fig:backward}
\end{figure}

Результатом обратного прохода является градиент по входу сети и набор
градиентов по слоям. Для линейного слоя градиент параметров
представлен парой \texttt{Matrix, Vector}: градиентом по матрице
весов и градиентом по вектору смещения. У слоя активации обучаемых
параметров нет, поэтому его градиент параметров пустой.

После этого \texttt{NeuralNetwork::update} идет по слоям слева направо
и передает каждому слою его градиенты, оптимизатор и локальное
состояние оптимизатора. Состояние оптимизатора отделено от кэша
прямого прохода. Кэш прямого прохода нужен только для одного шага
\texttt{forward -> backward}; состояние оптимизатора живет между
итерациями обучения и хранится у \texttt{Trainer}.

\subsection{Логика владения кэшами}

\texttt{std::any} используется для хранения объектов, тип которых
известен только конкретному компоненту. Для кэша прямого прохода
действует простое правило: объект, созданный слоем в \texttt{forward},
должен вернуться в тот же слой в \texttt{backward}. Сеть обеспечивает
порядок передачи, но не проверяет внутренний тип кэша.

Такой подход сохраняет локальность реализации. Если добавляется новый
слой, не требуется менять общий тип кэша или расширять общий
\texttt{std::variant}. Новый слой определяет собственный внутренний
тип кэша, возвращает его из \texttt{forward} и извлекает его в
\texttt{backward}.

Разделение кэшей также фиксирует границы ответственности. Слой
отвечает за локальную математику и локальный кэш.
\texttt{NeuralNetwork} отвечает за порядок обхода слоев.
\texttt{Trainer} связывает сеть, данные, функцию потерь и оптимизатор,
но не интерпретирует содержимое кэшей слоев.

\section{Базовые типы и генерация случайных величин}

\subsection{Базовые числовые типы}

Все вычисления в текущей версии выполняются в типе \texttt{double}.
Этот тип обеспечивает достаточную точность для реализованных
алгоритмов и не требует параметризации всего публичного API.

Для параметров, которые легко перепутать с обычными числами,
используется шаблон \texttt{StrongType}. На его основе определены типы
\texttt{LR}, \texttt{Gain}, \texttt{Beta1}, \texttt{Beta2},
\texttt{Eps}. Обертка хранит значение и имеет явное создание из
базового типа. Например, шаг обучения задается как
\texttt{LR\{0.001\}}, а не как обычный \texttt{double}. Это делает
вызовы конструкторов более читаемыми.

\subsection{RandomGenerator}

Класс \texttt{RandomGenerator} инкапсулирует \texttt{std::mt19937}. Он
создается с seed по умолчанию или с явно заданным seed:
\begin{cppcode}
nns::RandomGenerator rng(nns::Seed{42});
\end{cppcode}
Класс предоставляет методы \texttt{reseed}, \texttt{seed},
\texttt{init\_matrix}, \texttt{shuffle} и \texttt{get\_random\_value}.
Наиболее важным является метод \texttt{init\_matrix}, который
заполняет матрицу случайными значениями из нормального распределения.

Поддерживаются три схемы инициализации:
\begin{itemize}
    \item \texttt{Normal}: стандартное отклонение равно \texttt{gain};
    \item \texttt{XavierNormal}: стандартное отклонение равно
        $gain\sqrt{2/(fan_{in}+fan_{out})}$;
    \item \texttt{HeNormal}: стандартное отклонение равно
        $gain\sqrt{2/fan_{in}}$.
\end{itemize}
Инициализация Xavier предложена как способ стабилизировать дисперсию
сигналов при прохождении через глубокую
сеть~\cite{glorot2010understanding}. Инициализация He адаптирована для
сетей с rectifier-активациями, в частности с
ReLU~\cite{he2015delving}.

Метод \texttt{init\_matrix} проверяет, что матрица непустая, а
\texttt{gain} положителен. Проверка непустоты важна, поскольку формулы
инициализации используют размеры матрицы, а заполнение пустой матрицы
обычно свидетельствует об ошибке в конфигурации слоя.

\section{Функции потерь}

\subsection{Общий контракт}

Модуль функций потерь состоит из встроенных реализаций и type-erased
обертки \texttt{AnyLossFunction}. Интерфейс функции потерь описан в
разделе о стирании типов; здесь рассматриваются свойства встроенных
реализаций. Проверка формы вынесена во вспомогательную функцию
\texttt{validate\_same\_nonempty\_shape}. Все встроенные функции
потерь требуют, чтобы \texttt{y\_hat} и \texttt{y} имели одинаковую
форму и были непустыми. Функция потерь возвращает градиент той же
формы, что и выход сети, чтобы его можно было передать в
\texttt{NeuralNetwork::backward}.

\subsection{Loss-функции для регрессии}

Для регрессионных задач реализованы \texttt{MSELoss}, \texttt{MAELoss}
и \texttt{HuberLoss}. MSE чувствительна к большим ошибкам, поскольку
ошибка возводится в квадрат. MAE менее чувствительна к выбросам, но ее
градиент имеет разрыв в нуле. Huber loss объединяет свойства MSE и
MAE: при малой ошибке она квадратична, а при большой ошибке растет
линейно. Поэтому Huber loss часто используют как устойчивую
альтернативу MSE.

В реализации \texttt{HuberLoss} параметр \texttt{delta} является
публичным полем со значением по умолчанию $1.0$. Перед вычислением
значения и градиента проверяется, что \texttt{delta > 0}. Это пример
локальной валидации гиперпараметров: некорректное значение
обнаруживается в функции, для которой оно имеет смысл.

\subsection{Loss-функции для классификации}

Для бинарной классификации \texttt{BCELoss} работает с вероятностями.
Чтобы избежать вычисления $\log(0)$, предсказания ограничиваются снизу
и сверху через \texttt{eps}. Значение \texttt{eps} должно лежать в
интервале $(0, 0.5)$: при \texttt{eps <= 0} защита не работает, а при
\texttt{eps >= 0.5} отсечение разрушает смысл вероятностей.

\texttt{BCEWithLogitsLoss} принимает не вероятности, а logits. Это
предпочтительно, когда последним слоем модели является линейный слой
без сигмоиды. Объединение сигмоиды и бинарной кросс-энтропии в одной
формуле улучшает численную устойчивость и уменьшает риск переполнения.

Для многоклассовой классификации \texttt{CrossEntropyLoss} принимает
logits и one-hot разметку. Softmax применяется внутри функции потерь,
а градиент имеет простой вид:
\[
    \frac{\partial L}{\partial z} = \frac{p - y}{B},
\]
где $p$ --- результат softmax, $y$ --- one-hot вектор класса, $B$ ---
размер батча. Такое объединение softmax и кросс-энтропии соответствует
типичному устройству современных библиотек глубокого обучения.

\section{Оптимизаторы и scheduler-ы}

\subsection{Learning rate scheduler}

Scheduler отвечает только за значение шага обучения на текущей
итерации. В проекте реализованы \texttt{ConstantLR} и
\texttt{TimeDecayLR}. \texttt{ConstantLR} возвращает постоянное
значение, а \texttt{TimeDecayLR} вычисляет
\[
    \eta_t = \eta_0 \left(\frac{s_0}{s_0 + t}\right)^p.
\]
Здесь $\eta_0$ --- начальный шаг, $t$ --- номер итерации, $s_0$ и $p$
--- параметры убывания. В конструкторе можно задать любой из этих параметров, кроме $t$.

Оба scheduler-а хранят счетчик итераций. Оптимизатор вызывает
\texttt{iter\_step} после обработки каждого батча. Это означает, что
расписание шага обучения изменяется по итерациям оптимизации, а не по
эпохам.

\subsection{SGDOptimizer}

Стохастический градиентный спуск обновляет параметр по правилу
\[
    \theta_{t+1} = \theta_t - \eta_t \nabla_{\theta} L.
\]
В реализации \texttt{SGDOptimizer} один и тот же шаблонный приватный
метод \texttt{update\_any} применяется к матрицам и векторам. Перед
обновлением проверяется совпадение формы параметра и градиента.
Состояние оптимизатора для SGD отсутствует, поэтому метод
\texttt{update\_weights} возвращает пустой \texttt{std::any}.

\subsection{AdamOptimizer}

Adam использует экспоненциальные скользящие средние первого и второго
моментов градиента~\cite{kingma2014adam}. Для градиента $g_t$
обновления имеют вид
\[
    m_t = \beta_1 m_{t-1} + (1-\beta_1)g_t,
\]
\[
    v_t = \beta_2 v_{t-1} + (1-\beta_2)g_t^2.
\]
Затем применяется коррекция смещения:
\[
    \hat{m}_t = \frac{m_t}{1-\beta_1^t},
    \qquad
    \hat{v}_t = \frac{v_t}{1-\beta_2^t},
\]
и параметр обновляется по правилу
\[
    \theta_{t+1} = \theta_t - \eta_t
    \frac{\hat{m}_t}{\sqrt{\hat{v}_t}+\varepsilon}.
\]
В реализации состояние Adam хранится отдельно для каждого параметра в
кэше оптимизатора. Для матрицы весов и вектора смещения линейного слоя
создаются независимые кэши. При первом обновлении, когда кэш пуст,
моменты инициализируются нулями той же формы, что и градиент.

Значения \texttt{beta1}, \texttt{beta2} и \texttt{eps} проверяются
перед шагом Adam. Параметры \texttt{beta1} и \texttt{beta2} должны
лежать в диапазоне $[0,1)$, а \texttt{eps} должен быть положительным.
Нарушение этих ограничений приводит к исключению
\texttt{std::invalid\_argument}.

\section{Работа с данными}

\subsection{DataLoader}

Класс \texttt{DataLoader} хранит матрицы признаков и целевых значений,
размер батча, флаг перемешивания и вектор индексов. Конструктор
принимает матрицы по значению и перемещает их во внутренние поля:
\begin{cppcode}
nns::DataLoader loader(
    std::move(X),
    std::move(Y),
    nns::BatchSize{64},
    nns::Shuffle{true});
\end{cppcode}
При создании проверяется, что число объектов в \texttt{X} и \texttt{Y}
совпадает, размер батча положителен, а набор данных не пуст. Число
объектов определяется по числу столбцов. Метод \texttt{num\_batches}
возвращает только число полных батчей, то есть остаток в конце эпохи
отбрасывается.

Метод \texttt{get\_batch} создает новые матрицы батча и копирует в них
столбцы исходных матриц в порядке, заданном вектором индексов. Если
включено перемешивание, метод \texttt{reset\_epoch(RandomGenerator\&
rng)} перемешивает индексы. Вызов \texttt{reset\_epoch()} без
генератора допустим только при \texttt{Shuffle\{false\}}; иначе
библиотека не может выполнить воспроизводимое перемешивание и
выбрасывает \texttt{std::logic\_error}.

\subsection{CSV-загрузка}

Функция \texttt{load\_csv} загружает числовой CSV-файл и разделяет
колонки на признаки и целевые значения. Пользователь передает путь,
множество индексов целевых колонок, разделитель и флаг наличия
заголовка:
\begin{cppcode}
auto [X, Y] = nns::load_csv("data.csv", {0}, ',', true);
\end{cppcode}
Если первая колонка является меткой класса, а остальные колонки
являются признаками, вызов с \texttt{\{0\}} вернет матрицу признаков
без первой колонки и матрицу целей из первой колонки. Данные в
итоговых матрицах хранятся по принятому соглашению: объекты становятся
столбцами.

Функция выполняет несколько проверок. Множество целевых колонок не
должно быть пустым; файл должен открываться; файл должен содержать
хотя бы одну строку данных; все строки должны иметь одинаковое число
колонок; индексы целевых колонок должны попадать в диапазон; должна
существовать хотя бы одна feature-колонка. Эти проверки особенно важны
для CSV, поскольку ошибки формата данных иначе проявились бы позже как
ошибки размерности матриц.

\section{Trainer и цикл обучения}

\subsection{Роль Trainer}

Класс \texttt{Trainer} связывает четыре объекта: нейронную сеть,
оптимизатор, функцию потерь и загрузчик данных. Он хранит ссылки на
эти объекты, поэтому они должны жить дольше самого тренера. Такой
выбор предотвращает дорогие копирования сети и данных, а также делает
владение объектами явным на стороне пользователя.

Конструктор имеет вид:
\begin{cppcode}
nns::Trainer trainer(net, opt, loss, loader);
\end{cppcode}
Объект \texttt{Trainer} предоставляет методы \texttt{fit\_epoch},
\texttt{fit}, а также \texttt{reset\_optimizer\_state}. Методы
\texttt{fit\_epoch} обучают одну эпоху, а методы \texttt{fit} обучает сразу все эпохи
и возвращает вектор средних значений функции потерь по эпохам

\subsection{Алгоритм одной эпохи}

Одна эпоха обучения состоит из следующих шагов:
\begin{enumerate}
    \item сбросить эпоху в \texttt{DataLoader} и при необходимости
        перемешать индексы; \item для каждого полного батча получить
        матрицы \texttt{bX} и \texttt{bY}; \item выполнить
        \texttt{net.forward} и получить предсказание и кэш слоев;
        \item вычислить значение функции потерь; \item вычислить
        градиент функции потерь по предсказанию; \item выполнить
        \texttt{net.backward} и получить градиенты слоев; \item
        выполнить \texttt{net.update}, передав градиенты, оптимизатор
        и кэш оптимизатора; \item вызвать \texttt{opt.iter\_step};
        \item добавить значение loss к сумме по эпохе.
\end{enumerate}
После обработки всех батчей возвращается среднее значение функции
потерь по числу батчей. Если \texttt{DataLoader} не содержит ни одного
полного батча, тренер выбрасывает \texttt{std::logic\_error}.

\subsection{Состояние оптимизатора}

Состояние оптимизатора хранится внутри \texttt{Trainer} в поле
\texttt{std::any opt\_cache\_}. Для SGD состояние пустое, а для Adam
содержит моменты для каждого параметра каждого слоя. Такое размещение
состояния логично: состояние оптимизации относится не к самой модели
как математической функции, а к конкретному процессу обучения.

Метод \texttt{reset\_optimizer\_state} очищает кэш. Это полезно, если
одну и ту же сеть нужно продолжить обучать с заново инициализированным
оптимизатором или если пользователь меняет оптимизатор между
экспериментами.

\section{Контракты API и обработка ошибок}

Компоненты библиотеки проверяют входные параметры на границах
публичного API. Это касается размерностей матриц, положительности
гиперпараметров, согласованности числа батчей, наличия данных и
доступности внешних файлов. Проверки выполняются рядом с тем кодом,
который обладает достаточным контекстом для точного сообщения об
ошибке. Например, слой проверяет размерности своего входа, функция
потерь проверяет совпадение формы предсказания и целевой матрицы, а
\texttt{DataLoader} проверяет размер батча и число объектов.

Такое распределение проверок сохраняет независимость компонентов.
Линейный слой можно тестировать без тренера, функцию потерь --- без
сети, оптимизатор --- на одной матрице или одном векторе. Полный
интеграционный тест нужен только для проверки того, что компоненты
корректно соединяются в цикле обучения.

\subsection{Стратегия обработки ошибок}

В библиотеке используются исключения \texttt{std::invalid\_argument},
\texttt{std::logic\_error}, \texttt{std::out\_of\_range} и
\texttt{std::runtime\_error}. Их смысл различается:
\begin{itemize}
    \item \texttt{std::invalid\_argument} означает некорректное
        значение аргумента: неправильная форма матрицы,
        неположительный learning rate, неверный параметр функции
        потерь; \item \texttt{std::logic\_error} означает неправильный
        порядок действий: например, вызов \texttt{fit\_epoch()} без
        генератора при включенном перемешивании; \item
        \texttt{std::out\_of\_range} используется для обращения к
        несуществующему батчу; \item \texttt{std::runtime\_error}
        используется для ошибок внешней среды, например невозможности
        открыть CSV-файл.
\end{itemize}
Такое разделение делает сообщения об ошибках полезнее для
пользователя. Например, ошибка размера матрицы обычно исправляется в
коде конфигурации модели, а ошибка открытия CSV-файла может быть
связана с путями, рабочей директорией или отсутствием данных.

\section{Размерности, сложность и численная устойчивость}

\subsection{Соглашения о размерностях}

Большая часть ошибок при ручной реализации нейронных сетей связана с
перепутанными размерностями. В проекте принято соглашение, что объекты
выборки хранятся в столбцах. Поэтому линейный слой с входной
размерностью $d_{in}$ и выходной размерностью $d_{out}$ работает со
следующими объектами:
\begin{center}
\begin{tabular}{|c|c|c|}
\hline
Объект & Смысл & Размерность \\
\hline
$X$ & вход батча & $d_{in} \times B$ \\
$A$ & матрица весов & $d_{out} \times d_{in}$ \\
$b$ & вектор смещения & $d_{out} \times 1$ \\
$Y$ & выход батча & $d_{out} \times B$ \\
$dY$ & градиент по выходу & $d_{out} \times B$ \\
$dA$ & градиент по весам & $d_{out} \times d_{in}$ \\
$db$ & градиент по смещению & $d_{out} \times 1$ \\
$dX$ & градиент по входу & $d_{in} \times B$ \\
\hline
\end{tabular}
\end{center}
Именно эти размерности проверяются в коде. Например,
\texttt{LinearLayer::predict} требует, чтобы число строк входной
матрицы совпадало с числом столбцов матрицы весов.
\texttt{LinearLayer::backward} проверяет, что число строк \texttt{dY}
совпадает с выходной размерностью, а число столбцов совпадает с
размером батча, сохраненным в кэше.

\subsection{Асимптотическая сложность}

Основная вычислительная стоимость полносвязной сети приходится на
матричные умножения. Для одного линейного слоя прямой проход требует
умножения матрицы $d_{out} \times d_{in}$ на матрицу $d_{in} \times
B$, что имеет сложность
\[
    O(d_{out} d_{in} B).
\]
Прибавление смещения имеет сложность $O(d_{out}B)$ и обычно
существенно дешевле умножения.

Обратный проход линейного слоя содержит два матричных умножения:
вычисление $dA = dY X^T$ и вычисление $dX = A^T dY$. Каждое из них
также имеет порядок $O(d_{out} d_{in} B)$. Следовательно, обучение
одного линейного слоя на одном батче имеет ту же асимптотику, что и
несколько матричных умножений с теми же размерами. Для сети из $m$
линейных слоев суммарная стоимость одной итерации равна сумме
стоимостей по слоям.

Слой активации имеет сложность $O(n)$, где $n$ --- число элементов
входной матрицы. Он дешевле линейного слоя, если размерности
достаточно велики. Функции потерь также обычно имеют линейную
сложность по числу элементов выхода. Исключением не является и
softmax: он требует прохода по каждому столбцу для поиска максимума,
вычисления экспонент и нормировки, что дает $O(kB)$ для $k$ классов и
размера батча $B$.

\subsection{Память и кэши}

Логика владения кэшами описана в архитектурной части отчета. С точки
зрения памяти важно, что во время прямого прохода сохраняются данные,
необходимые для последующего вычисления локальных градиентов. Поэтому
память кэша прямого прохода имеет тот же порядок, что и сумма размеров
входов всех слоев на данном батче.

Состояние Adam также требует дополнительной памяти. Для каждого
параметра хранятся две величины той же формы: первый и второй моменты.
Для линейного слоя это две матрицы размера $d_{out} \times d_{in}$ и
два вектора размера $d_{out}$. Поэтому Adam примерно утраивает память,
связанную с параметрами: сами параметры плюс два набора моментов. SGD
такого состояния не имеет.

\subsection{Численная устойчивость}

В нескольких местах реализации явно учитывается численная
устойчивость. В \texttt{CrossEntropyLoss::softmax} из каждого столбца
logits вычитается максимум. Без этого при больших положительных logits
экспонента может переполниться. Так как softmax инвариантен к
прибавлению одной и той же константы ко всем компонентам вектора,
такое преобразование не меняет математический результат.

В \texttt{BCELoss} вероятности ограничиваются через \texttt{eps}. Это
защищает от логарифма нуля и деления на ноль в градиенте. В
\texttt{BCEWithLogitsLoss} используется стабильная формула через
$\max(z,0)$ и $\log(1+e^{-|z|})$, которая избегает прямого вычисления
$\log(1+e^z)$ при большом $z$.

В Adam параметр $\varepsilon$ добавляется к знаменателю при делении на
корень из второго момента. Это предотвращает деление на ноль и
стабилизирует обновление при очень малых градиентах.

\section{Пример пользовательского сценария}

\subsection{Минимальный пример обучения}

Ниже приведен сокращенный пример обучения сети на небольшом наборе
данных. Он показывает полный путь от создания данных до запуска
\texttt{Trainer}.
\begin{cppcode}
#include <iostream>
#include <nns/NNS.hpp>

int main() {
    nns::RandomGenerator rng(nns::Seed{42});

    nns::Matrix X(2, 4);
    X << 0.0, 0.0, 1.0, 1.0,
         0.0, 1.0, 0.0, 1.0;

    nns::Matrix Y = nns::Matrix::Zero(2, 4);
    Y(0, 0) = 1.0;
    Y(1, 1) = 1.0;
    Y(1, 2) = 1.0;
    Y(0, 3) = 1.0;

    nns::DataLoader loader(std::move(X), std::move(Y),
                           nns::BatchSize{4}, nns::Shuffle{true});

    nns::NeuralNetwork net(
        nns::LinearLayer(
            nns::In{2},
            nns::Out{8},
            rng,
            nns::Distribution::HeNormal),
        nns::ReLU(),
        nns::LinearLayer(
            nns::In{8},
            nns::Out{2},
            rng,
            nns::Distribution::XavierNormal));

    nns::AnyOptimizer opt =
        nns::make_AnyOptimizer(
            nns::AdamOptimizer(
                nns::ConstantLR(nns::LR{0.01})));
    nns::AnyLossFunction loss =
        nns::make_AnyLossFunction<nns::CrossEntropyLoss>();

    nns::Trainer trainer(net, opt, loss, loader);
    std::vector<nns::Scalar> history = trainer.fit(200, rng);

    std::cout << "last loss = " << history.back() << '\n';
}
\end{cppcode}
В этом примере используется \texttt{CrossEntropyLoss}, поэтому
последний слой возвращает logits, а softmax применяется внутри функции
потерь. Для инференса вероятности можно получить вызовом
\texttt{CrossEntropyLoss::softmax} от результата \texttt{net.predict}.

\section{Экспериментальное обучение на MNIST}

\subsection{Конфигурация эксперимента}

В файле \texttt{main.cpp} реализован эксперимент на наборе MNIST.
MNIST содержит изображения рукописных цифр размера $28 \times 28$ и
разметку по десяти классам~\cite{lecun1998gradient}. В репозитории
эксперимент рассчитан на CSV-представление данных: первая колонка
содержит метку класса, остальные 784 колонки содержат значения
пикселей.

Обработка данных выполняется следующим образом. Функция
\texttt{load\_csv} отделяет колонку метки от признаков. Значения
пикселей нормируются делением на 255, то есть переводятся из диапазона
$[0,255]$ в диапазон $[0,1]$. Метки классов преобразуются в one-hot
представление размерности $10 \times N$. После этого обучающая выборка
передается в \texttt{DataLoader} с размером батча 100 и включенным
перемешиванием.

Схема эксперимента показана на рисунке~\ref{fig:mnist}. Архитектура
сети:
\[
    784 \rightarrow 128 \rightarrow 64 \rightarrow 10.
\]
Между линейными слоями используются \texttt{ReLU}. Последний слой
возвращает logits. В качестве функции потерь используется
\texttt{CrossEntropyLoss}, которая применяет softmax внутри вычисления
loss и gradient. Оптимизатор --- Adam с постоянным learning rate
$0.001$. Число эпох равно 10, seed генератора равен 42.

\begin{figure}[!t]
    \centering
    \includesvg[inkscapelatex=false,width=1\linewidth]{mnist_pipeline}
    \caption{Поток данных в эксперименте MNIST}
    \label{fig:mnist}
\end{figure}

Фрагмент кода, задающий конфигурацию обучения, приведен ниже.
\begin{cppcode}
RandomGenerator rng(Seed{42});

constexpr size_t batch_size = 100;
constexpr size_t epochs = 10;

DataLoader train_loader(
    std::move(train_X),
    std::move(train_Y),
    BatchSize{batch_size},
    Shuffle{true});

NeuralNetwork net(
    LinearLayer(
        In{784},
        Out{128},
        rng,
        Distribution::Normal,
        Gain{0.01}),
    ReLU(),
    LinearLayer(
        In{128},
        Out{64},
        rng,
        Distribution::Normal,
        Gain{0.01}),
    ReLU(),
    LinearLayer(
        In{64},
        Out{10},
        rng,
        Distribution::Normal,
        Gain{0.01}));

AnyOptimizer opt = make_AnyOptimizer(
    AdamOptimizer(ConstantLR(LR{0.001})));

AnyLossFunction loss_func = make_AnyLossFunction<CrossEntropyLoss>();
Trainer trainer(net, opt, loss_func, train_loader);
\end{cppcode}

\subsection{Методика расчета метрик}

Для оценки качества используется тестовая часть MNIST. Предсказанный
класс определяется как индекс максимального logit:
\[
    \hat{c} = \arg\max_i z_i.
\]
На основе истинного класса $c$ и предсказанного класса $\hat{c}$
строится матрица ошибок $C$, где $C_{ij}$ равно числу объектов
истинного класса $i$, отнесенных моделью к классу $j$.

Accuracy вычисляется как доля правильно классифицированных объектов:
\[
    accuracy = \frac{\sum_i C_{ii}}{\sum_{i,j} C_{ij}}.
\]
Для каждого класса дополнительно вычисляются precision, recall и F1:
\[
    precision_i = \frac{TP_i}{TP_i + FP_i}, \qquad
    recall_i = \frac{TP_i}{TP_i + FN_i},
\]
\[
    F1_i =
    \frac{2 \cdot precision_i \cdot recall_i}
         {precision_i + recall_i}.
\]
Macro-метрики являются средним арифметическим по классам,
weighted-метрики являются средним с весами, равными числу объектов
соответствующего класса.

Код расчета матрицы ошибок и accuracy использует тот же
\texttt{predict}, который применяется в пользовательском режиме
инференса:
\begin{cppcode}
Matrix logits = net.predict(batch_X);
for (Index col = 0; col < cols; ++col) {
    Index pred_class = 0;
    Index true_class = 0;
    logits.col(col).maxCoeff(&pred_class);
    batch_Y.col(col).maxCoeff(&true_class);
    confusion[true_class][pred_class]++;
    if (pred_class == true_class) {
        ++correct;
    }
}
\end{cppcode}

\subsection{Результаты обучения}

Эксперимент запускался в конфигурации \texttt{release}. В
debug-конфигурации тот же код выполняется существенно медленнее из-за
отсутствия оптимизаций компилятора. История функции потерь на
обучающей выборке и accuracy на тестовой выборке:
\begin{center}
\begin{tabular}{|c|c|c|}
\hline
Эпоха & Train loss & Test accuracy \\
\hline
1 & 0.534551 & 0.9197 \\
2 & 0.235347 & 0.9438 \\
3 & 0.164380 & 0.9538 \\
4 & 0.125233 & 0.9631 \\
5 & 0.102050 & 0.9653 \\
6 & 0.082142 & 0.9717 \\
7 & 0.067856 & 0.9734 \\
8 & 0.057028 & 0.9723 \\
9 & 0.047776 & 0.9754 \\
10 & 0.039908 & 0.9751 \\
\hline
\end{tabular}
\end{center}

Итоговые метрики на тестовой выборке после 10 эпох:
\begin{center}
\begin{tabular}{|l|c|}
\hline
Метрика & Значение \\
\hline
Test cross-entropy & 0.084976 \\
Accuracy & 0.975100 \\
Macro precision & 0.975139 \\
Macro recall & 0.974866 \\
Macro F1 & 0.974938 \\
Weighted precision & 0.975243 \\
Weighted recall & 0.975100 \\
Weighted F1 & 0.975108 \\
\hline
\end{tabular}
\end{center}

Матрица ошибок для тестовой выборки приведена ниже. Строки
соответствуют истинному классу, столбцы --- предсказанному классу.
\begin{center}
\small
\begin{tabular}{c|rrrrrrrrrr}
 & 0 & 1 & 2 & 3 & 4 & 5 & 6 & 7 & 8 & 9 \\
\hline
0 & 967 & 0 & 2 & 2 & 3 & 0 & 2 & 1 & 2 & 1 \\
1 & 0 & 1122 & 1 & 1 & 0 & 0 & 2 & 1 & 7 & 1 \\
2 & 3 & 4 & 997 & 8 & 5 & 1 & 1 & 7 & 6 & 0 \\
3 & 0 & 0 & 1 & 996 & 2 & 3 & 0 & 4 & 1 & 3 \\
4 & 1 & 0 & 0 & 0 & 962 & 0 & 4 & 4 & 1 & 10 \\
5 & 2 & 0 & 0 & 13 & 2 & 867 & 4 & 1 & 2 & 1 \\
6 & 5 & 2 & 1 & 0 & 15 & 8 & 926 & 0 & 1 & 0 \\
7 & 1 & 5 & 7 & 4 & 1 & 0 & 0 & 1004 & 0 & 6 \\
8 & 3 & 0 & 3 & 9 & 5 & 3 & 4 & 5 & 936 & 6 \\
9 & 1 & 2 & 0 & 7 & 13 & 2 & 0 & 8 & 2 & 974 \\
\end{tabular}
\end{center}

Полученные значения показывают, что реализованный обучающий цикл
корректно снижает функцию потерь и формирует пригодную для
классификации модель. Наиболее частые ошибки связаны с визуально
близкими цифрами, например с переходами между классами 4 и 9, 5 и 3, 7
и 2.

\section{Сборка проекта}

\subsection{Сборка через CMake}

Проект собирается через CMake. В корневом \texttt{CMakeLists.txt}
объявляется проект \texttt{NNS}, подключаются зависимости через
\texttt{FetchContent}, создается \texttt{INTERFACE}-цель \texttt{nns}
и псевдоним \texttt{NNS::nns}. Для библиотеки задается требование
\texttt{cxx\_std\_23}.

В проекте предусмотрены опции:
\begin{itemize}
    \item \texttt{NNS\_BUILD\_EXAMPLE} --- собирать пример из
        \texttt{main.cpp}; \item \texttt{NNS\_BUILD\_TESTING} ---
        собирать тесты; \item \texttt{NNS\_ENABLE\_FORMAT\_TARGETS}
        --- включить цели форматирования.
\end{itemize}
Файл \texttt{CMakePresets.json} определяет пресеты \texttt{debug} и
\texttt{release}. Оба используют Ninja, включают тесты и экспортируют
\texttt{compile\_commands.json}. Для удобства добавлен
\texttt{Justfile} с командами \texttt{configure}, \texttt{build},
\texttt{test}, \texttt{format}, \texttt{format-check} и
\texttt{check}.

\subsection{Сборка через just}
Так же поддерживается сборка через just. Для сборки в debug режиме используются следующие команды:
\begin{cppcode}
just configure
just build
just test
\end{cppcode}
Для прогона по тестам и применение clang-format используется:
\begin{cppcode}
just check
\end{cppcode}
Для release сборки к командам из debug режима добавляется флаг release:
\begin{cppcode}
just configure release
just build release
just test release
\end{cppcode}
\section{Тестирование}

\subsection{Организация тестового набора}

Тесты разделены по функциональным областям:
\begin{itemize}
    \item \texttt{activation\_tests.cpp} проверяет встроенные и
        пользовательские активации; \item \texttt{core\_tests.cpp}
        проверяет базовые типы; \item
        \texttt{layer\_network\_trainer\_tests.cpp} проверяет слои,
        сеть и тренер; \item \texttt{loss\_tests.cpp} проверяет
        функции потерь; \item \texttt{optimizer\_tests.cpp} проверяет
        scheduler-ы и оптимизаторы; \item \texttt{utils\_tests.cpp}
        проверяет генератор случайных чисел, \texttt{DataLoader} и
        CSV-загрузку.
\end{itemize}

\subsection{Проверяемые сценарии}

Тестовый набор покрывает три группы сценариев. Первая группа ---
проверка математических значений на малых входах. Например, для
\texttt{MSELoss}, \texttt{MAELoss}, \texttt{HuberLoss},
\texttt{BCELoss} и \texttt{CrossEntropyLoss} проверяются значения
функции потерь и отдельные элементы градиента. Для \texttt{ReLU},
\texttt{Sigmoid} и \texttt{Tanh} проверяются значения функции и
производной.

Вторая группа --- проверка контрактов размерностей. Тесты вызывают
слои, сеть, функции потерь и оптимизаторы с некорректными формами
матриц и ожидают исключения \texttt{std::invalid\_argument}. Например,
\texttt{LinearLayer::predict} должен отклонять вход с неправильным
числом строк, \texttt{LinearLayer::backward} должен отклонять градиент
неправильной формы, а \texttt{NeuralNetwork::backward} должен
отклонять кэш с числом элементов, не совпадающим с числом слоев.

Третья группа --- интеграционные сценарии. В них проверяется, что
\texttt{NeuralNetwork} поддерживает полный цикл \texttt{predict ->
forward -> backward -> update}, а \texttt{Trainer} выполняет эпоху
обучения, возвращает историю loss, поддерживает перемешивание через
\texttt{RandomGenerator} и отклоняет \texttt{DataLoader}, не
содержащий полного батча.

\subsection{Запуск тестов}

Тесты запускаются через CTest:
\begin{cppcode}
cmake --preset debug
cmake --build --preset debug
ctest --preset debug
\end{cppcode}
Также предусмотрена команда \texttt{just check}, которая собирает
проект, запускает тесты и проверяет форматирование:
\begin{cppcode}
just check
\end{cppcode}
На текущей версии проекта команда \texttt{ctest --preset debug}
выполняет 29 тестов; все 29 тестов завершаются успешно.

\subsection{Статические соглашения о стиле}

Для форматирования используется \texttt{clang-format}. Цели
\texttt{format} и \texttt{format-check} создаются CMake при наличии
исполняемого файла \texttt{clang-format}. В репозитории также
присутствует конфигурация \texttt{.clang-tidy}, ориентированная на
группы проверок \texttt{google-*}, \texttt{cppcoreguidelines-*} и
\texttt{readability-*}. Часть проверок, связанных с magic numbers,
отключена, поскольку в тестах и численных формулах небольшие константы
часто являются частью математического определения.

\section{Текущие ограничения}

\subsection{Ограничения модели}

Текущая версия библиотеки поддерживает только последовательные
полносвязные сети. В ней отсутствуют сверточные слои, рекуррентные
слои, residual-соединения, ветвления вычислительного графа и
автоматическое дифференцирование произвольных операций. Это осознанное
ограничение: проект фокусируется на понятной реализации базового
обучения многослойного перцептрона.

Функции активации являются только скалярными. Поэтому операции,
зависящие от всего вектора, например softmax как отдельный слой, не
входят в интерфейс \texttt{ActivationLayer}. В текущей реализации
softmax расположен внутри \texttt{CrossEntropyLoss}. Для задач, где
softmax нужен как самостоятельная операция инференса, пользователь
может вызвать \texttt{CrossEntropyLoss::softmax} явно.

\subsection{Ограничения производительности}

Библиотека ориентирована на CPU и использует Eigen. В ней нет
поддержки GPU, смешанной точности, распараллеливания на уровне батчей
или специализированных BLAS-бэкендов. Кроме того, кэши и батчи создают
дополнительные копии матриц. Для больших моделей потребуется
оптимизация владения памятью и уменьшение числа временных объектов.

% \section{Направления дальнейшего развития}

% Первое направление развития --- расширение набора слоев. Наиболее
% естественно добавить dropout, batch normalization и softmax-слой с
% матричным интерфейсом. Затем можно рассмотреть сверточные слои, однако
% они потребуют отдельного проектирования форматов тензоров и более
% сложных операций, чем текущие матрицы.

% Второе направление --- улучшение работы с данными. Полезно добавить
% поддержку неполного последнего батча, разбиение на train/validation
% выборки, нормализацию признаков, сохранение статистик нормализации и
% более устойчивую загрузку CSV. Также можно добавить сохранение и
% загрузку параметров модели.

% Третье направление --- развитие оптимизации. Помимо SGD и Adam можно
% реализовать momentum SGD, RMSProp, AdamW и weight decay. Для
% scheduler-ов полезны cosine decay, step decay и warmup. При этом
% текущая архитектура уже позволяет добавлять новые оптимизаторы без
% изменения \texttt{Trainer}, если новый оптимизатор удовлетворяет
% контракту \texttt{AnyOptimizer}.

% Четвертое направление --- улучшение численной и инженерной надежности.
% Для функций потерь можно добавить больше тестов на численные
% градиенты. Для слоев можно добавить gradient checking через конечные
% разности. Для сборки можно включить sanitizer-конфигурации и
% CI-сценарий, запускающий сборку, тесты и проверку форматирования.

\section{Заключение}

В рамках проекта разработана библиотека \texttt{NNS} для построения и
обучения полносвязных нейронных сетей на C++23. Реализованы генератор случайных чисел, линейные слои, слои активации,
последовательная нейронная сеть, функции потерь, scheduler-ы шага
обучения, оптимизаторы SGD и Adam, загрузчик данных, CSV-парсер и
обучающий цикл.

Ключевым архитектурным решением стало использование type erasure через
библиотеку \texttt{proxy}. Благодаря этому пользователь может
расширять библиотеку собственными функциями активации, функциями
потерь, оптимизаторами и scheduler-ами без наследования от
библиотечных базовых классов. Вторым важным решением стало явное
возвращение кэша прямого прохода и состояния оптимизатора через
\texttt{std::any}; это делает поток данных между forward, backward и
update прозрачным и локализует состояние обучения в \texttt{Trainer}.

Реализация покрыта модульными тестами, которые проверяют корректные
сценарии и ошибки нарушения контрактов. Проект собирается через CMake,
поддерживает пресеты debug/release, форматирование через clang-format
и подключение как \texttt{NNS::nns}. В результате реализован полный
стек обучения полносвязной нейронной сети: от загрузки и батчинга
данных до вычисления градиентов и обновления параметров оптимизатором.

\bibliographystyle{plainurl}
\bibliography{bibl}

\end{document}
